\title{Keyframe-based Visual-Inertial Odometry for Small Workspace}
\author{Xi Li}

% You can use the \\ command to list multiple previous degrees
%       \prevdegrees{B.S., University of California (1978) \\
%                    S.M., Massachusetts Institute of Technology (1981)}
\department{Department of Computer Science}

% If the thesis is for two degrees simultaneously, list them both
% separated by \and like this:
% \degree{Doctor of Philosophy \and Master of Science}
\degree{Master of Science in Computer Science}

% As of the 2007-08 academic year, valid degree months are September, 
% February, or June.  The default is June.
\degreemonth{June}
\degreeyear{2016}
\thesisdate{June 1, 2016}

%% By default, the thesis will be copyrighted to SAAR.  If you need to copyright
%% the thesis to yourself, just specify the `vi' documentclass option.  If for
%% some reason you want to exactly specify the copyright notice text, you can
%% use the \copyrightnoticetext command.  
% \copyrightnoticetext{\copyright Quynh Nguyen Ngoc. All rights reserved.}

% If there is more than one supervisor, use the \supervisor command
% once for each.
\supervisor{Roland Angst}{Dr.}

% This is the department committee chairman, not the thesis committee
% chairman.  You should replace this with your Department's Committee
% Chairman.
\chairman{Prof}{Chairman, Department Committee on Graduate Theses}

% Make the titlepage based on the above information.  If you need
% something special and can't use the standard form, you can specify
% the exact text of the titlepage yourself.  Put it in a titlepage
% environment and leave blank lines where you want vertical space.
% The spaces will be adjusted to fill the entire page.  The dotted
% lines for the signatures are made with the \signature command.

% \maketitle
\noindent
\begin{titlepage}
\includegraphics[width=4cm,natwidth=520,natheight=210]{logo_saarland.png}
\hfill
\begin{minipage}{8cm}
\centering
\vspace{-1.5cm} 
\textbf{Universit\"at des Saarlandes \\ Max-Planck-Institut f\"ur Informatik}
\end{minipage}
\hfill
\includegraphics[width=2.8cm,natwidth=300,natheight=244]{mpg.png}
\vfill

\Large
\textbf{Keyframe-based Visual-Inertial Odometry for Small Workspace}
\vfill

\large
Masterarbeit im Fach Informatik \\
Master's Thesis in Computer Science \\
von / by \\
\textbf{Xi Li} \\
\vfill

angefertigt unter der Leitung von / supervised by \\
\large{\textsc{Dr. Roland Angst}}\\
\vfill

betreut von / advised by \\
\large{\textsc{Dr. Roland Angst}} \\
\vfill

begutachtet von / reviewers \\
\large{\textsc{Dr. Roland Angst \\ Prof. Dr. Antonio Kr{\"u}ger}} \\
\vfill

Saarbr\"ucken, June 2016
\end{titlepage}


% The abstractpage environment sets up everything on the page except
% the text itself.  The title and other header material are put at the
% top of the page, and the supervisors are listed at the bottom.  A
% new page is begun both before and after.  Of course, an abstract may
% be more than one page itself.  If you need more control over the
% format of the page, you can use the abstract environment, which puts
% the word "Abstract" at the beginning and single spaces its text.

%% You can either \input (*not* \include) your abstract file, or you can put
%% the text of the abstract directly between the \begin{abstractpage} and
%% \end{abstractpage} commands.

% First copy: start a new page, and save the page number.
\cleardoublepage
% Uncomment the next line if you do NOT want a page number on your
% abstract and acknowledgments pages.
% \pagestyle{empty}
% \setcounter{savepage}{\thepage}
% \begin{abstractpage}
%     %% The text of your abstract and nothing else (other than comments) goes here.
%% It will be single-spaced and the rest of the text that is supposed to go on
%% the abstract page will be generated by the abstractpage environment.  This
%% file should be \input (not \include 'd) from cover.tex.

Mobile robot navigation has been an attractive topic in the Robotics area for a long time. \textit{Visual SLAM} simultaneously locates robots' poses and constructs a map via analysing measurements obtained from visual sensors. However due to the costly computational complexity of camera, it is non-trivial to gain a high-quality pose estimation while the whole system runs in real-time. \textit{Inertial Measurement Unit} (IMU) measures the rotational rate and acceleration of the robot under the higher output frequency, and it has been considered as an appropriate compensation to a single camera in such navigation systems. In this master thesis, we aim to exploit a visual-inertial odometry, to further improve localization results by fusing visual data and IMU data. More precisely, we apply a loosely-coupled visual-inertial integration method to keep our system runs in a constant time. 

We first discuss the types of visual SLAM. Generally, visual SLAM has two traditional approaches, which are filter-based methods and keyframe-based methods respectively. Filter-based methods marginalise all the previous camera poses to obtain the current pose estimation as well as optimized landmarks, whereas keyframe-based methods keep some of former poses and applies bundle adjustment to further improve the estimation results. We finally choose keyframe-based visual SLAM to be correction data in our Kalman filter framework because it is more efficient when scale of scene becomes larger, another reason is that it can handle loop closures more easily. We then discuss the way to integrate IMU data. We first learn \textit{quaternion} and its operations. Instead of using Euler angle (Gimbal Lock) or rotational matrix (high memory cost), we formulate the system rotation parameters with quaternion. Followed by that, we explain the time-integration and differentiations related to the quaternions, and we suggest a error-state Kalman filter to propagate system states while tracking the system uncertainties. In order to fuse the visual data and IMU data, we suggest an adapted map scale method to automatically propagate the map scale factor, therefore solve the unknown map scale problem existed in mono visual SLAM. A keyframe-based bundle adjustment is then exploited for improving the tracking quality. Our experiments show that our suggested framework obtain more accurate and stable results compared to the state-of-art visual odometry, and our system can be easily extended to large scaled scene.

In addition to Kalman filter based framework, we briefly introduce a manifold optimization approach to solve visual-inertial odometry problem. We have suggested a standard way to solve manifold optimization as well as several ideas to reduce the computational time in the end.


% \end{abstractpage}

% Additional copy: start a new page, and reset the page number.  This way,
% the second copy of the abstract is not counted as separate pages.
% Uncomment the next 6 lines if you need two copies of the abstract
% page.
% \setcounter{page}{\thesavepage}
% \begin{abstractpage}
% %% The text of your abstract and nothing else (other than comments) goes here.
%% It will be single-spaced and the rest of the text that is supposed to go on
%% the abstract page will be generated by the abstractpage environment.  This
%% file should be \input (not \include 'd) from cover.tex.

Mobile robot navigation has been an attractive topic in the Robotics area for a long time. \textit{Visual SLAM} simultaneously locates robots' poses and constructs a map via analysing measurements obtained from visual sensors. However due to the costly computational complexity of camera, it is non-trivial to gain a high-quality pose estimation while the whole system runs in real-time. \textit{Inertial Measurement Unit} (IMU) measures the rotational rate and acceleration of the robot under the higher output frequency, and it has been considered as an appropriate compensation to a single camera in such navigation systems. In this master thesis, we aim to exploit a visual-inertial odometry, to further improve localization results by fusing visual data and IMU data. More precisely, we apply a loosely-coupled visual-inertial integration method to keep our system runs in a constant time. 

We first discuss the types of visual SLAM. Generally, visual SLAM has two traditional approaches, which are filter-based methods and keyframe-based methods respectively. Filter-based methods marginalise all the previous camera poses to obtain the current pose estimation as well as optimized landmarks, whereas keyframe-based methods keep some of former poses and applies bundle adjustment to further improve the estimation results. We finally choose keyframe-based visual SLAM to be correction data in our Kalman filter framework because it is more efficient when scale of scene becomes larger, another reason is that it can handle loop closures more easily. We then discuss the way to integrate IMU data. We first learn \textit{quaternion} and its operations. Instead of using Euler angle (Gimbal Lock) or rotational matrix (high memory cost), we formulate the system rotation parameters with quaternion. Followed by that, we explain the time-integration and differentiations related to the quaternions, and we suggest a error-state Kalman filter to propagate system states while tracking the system uncertainties. In order to fuse the visual data and IMU data, we suggest an adapted map scale method to automatically propagate the map scale factor, therefore solve the unknown map scale problem existed in mono visual SLAM. A keyframe-based bundle adjustment is then exploited for improving the tracking quality. Our experiments show that our suggested framework obtain more accurate and stable results compared to the state-of-art visual odometry, and our system can be easily extended to large scaled scene.

In addition to Kalman filter based framework, we briefly introduce a manifold optimization approach to solve visual-inertial odometry problem. We have suggested a standard way to solve manifold optimization as well as several ideas to reduce the computational time in the end.


% \end{abstractpage}

\cleardoublepage
\pagestyle{empty}
\setcounter{savepage}{\thepage}
\subsubsection*{Eidesstattliche Erkl{\"a}rung}
Ich erkl{\"a}re hiermit an Eides Statt, dass ich die vorliegende Arbeit selbstst{\"a}ndig verfasst und
keine anderen als die angegebenen Quellen und Hilfsmittel verwendet habe.
\subsubsection*{Statement in Lieu of an Oath}
I hereby confirm under oath that I have written this thesis on my own and that I have not
used any other media or materials than the ones referred to in this thesis.

\vspace{1cm}

\subsubsection*{Eidesstattliche Erkl{\"a}rung}
Ich erkl{\"a}re hiermit an Eides Statt, dass die vorliegende Arbeit mit der elektronischen Version {\"u}bereinstimmt.
\subsubsection*{Statement in Lieu of an Oath}
I hereby confirm the congruence of the contents of the printed data and the electronic version of the thesis.

\vspace{1cm}

\subsubsection*{Einverst{\"a}ndniserkl{\"a}rung}
Ich bin damit einverstanden, dass meine (bestandene) Arbeit in beiden Versionen in die
Bibliothek der Informatik aufgenommen und damit ver{\"o}ffentlicht wird.

\subsubsection*{Declaration of Consent}
I agree to make both versions of my thesis (with a passing grade) accessible to the public 
by having them added to the library of the Computer Science Department.

\vspace{3cm}

\par\noindent\makebox[2in]{\hrulefill} 		\hfill\makebox[2in]{\hrulefill}%
\par\noindent\makebox[2in][l]{(Datum / Date)}	\hfill\makebox[2in][l]{(Unterschrift / Signature)}%

\newpage
\section*{Acknowledgments}


\pagestyle{empty}
\setcounter{savepage}{\thepage}
\begin{abstractpage}
    %% The text of your abstract and nothing else (other than comments) goes here.
%% It will be single-spaced and the rest of the text that is supposed to go on
%% the abstract page will be generated by the abstractpage environment.  This
%% file should be \input (not \include 'd) from cover.tex.

Mobile robot navigation has been an attractive topic in the Robotics area for a long time. \textit{Visual SLAM} simultaneously locates robots' poses and constructs a map via analysing measurements obtained from visual sensors. However due to the costly computational complexity of camera, it is non-trivial to gain a high-quality pose estimation while the whole system runs in real-time. \textit{Inertial Measurement Unit} (IMU) measures the rotational rate and acceleration of the robot under the higher output frequency, and it has been considered as an appropriate compensation to a single camera in such navigation systems. In this master thesis, we aim to exploit a visual-inertial odometry, to further improve localization results by fusing visual data and IMU data. More precisely, we apply a loosely-coupled visual-inertial integration method to keep our system runs in a constant time. 

We first discuss the types of visual SLAM. Generally, visual SLAM has two traditional approaches, which are filter-based methods and keyframe-based methods respectively. Filter-based methods marginalise all the previous camera poses to obtain the current pose estimation as well as optimized landmarks, whereas keyframe-based methods keep some of former poses and applies bundle adjustment to further improve the estimation results. We finally choose keyframe-based visual SLAM to be correction data in our Kalman filter framework because it is more efficient when scale of scene becomes larger, another reason is that it can handle loop closures more easily. We then discuss the way to integrate IMU data. We first learn \textit{quaternion} and its operations. Instead of using Euler angle (Gimbal Lock) or rotational matrix (high memory cost), we formulate the system rotation parameters with quaternion. Followed by that, we explain the time-integration and differentiations related to the quaternions, and we suggest a error-state Kalman filter to propagate system states while tracking the system uncertainties. In order to fuse the visual data and IMU data, we suggest an adapted map scale method to automatically propagate the map scale factor, therefore solve the unknown map scale problem existed in mono visual SLAM. A keyframe-based bundle adjustment is then exploited for improving the tracking quality. Our experiments show that our suggested framework obtain more accurate and stable results compared to the state-of-art visual odometry, and our system can be easily extended to large scaled scene.

In addition to Kalman filter based framework, we briefly introduce a manifold optimization approach to solve visual-inertial odometry problem. We have suggested a standard way to solve manifold optimization as well as several ideas to reduce the computational time in the end.


\end{abstractpage}

%%%%%%%%%%%%%%%%%%%%%%%%%%%%%%%%%%%%%%%%%%%%%%%%%%%%%%%%%%%%%%%%%%%%%%
% -*-latex-*-
